\documentclass[11pt,oneside]{article}
\usepackage[a4paper,margin=27mm,headheight=15pt]{geometry}
\usepackage[T1]{fontenc}
\usepackage[utf8]{inputenc}
\IfFileExists{libertinus.sty}{\usepackage{libertinus}}{\usepackage{lmodern}}
\usepackage{microtype}
\usepackage{amsmath,amssymb,amsthm,mathtools,mathrsfs}
\usepackage{bm}
\usepackage{booktabs,longtable,array,tabularx}
\usepackage{graphicx}
\usepackage{pgfplots}
\usepackage{xcolor}
\usepackage{enumitem}
\usepackage{fancyhdr}
\usepackage{titlesec}
\usepackage{listings}
\usepackage{xurl}
\usepackage{hyperref}
\usepackage[nameinlink,noabbrev]{cleveref}
\definecolor{linkblue}{RGB}{25,75,135}
\definecolor{shadegray}{RGB}{246,247,249}
\definecolor{rulegray}{RGB}{195,201,208}
\pgfplotsset{compat=1.18}
\hypersetup{
  colorlinks=true,
  linkcolor=linkblue,
  citecolor=linkblue,
  urlcolor=linkblue,
  pdftitle={The Dyadic Web of the Fabius Function},
  pdfsubject={Moment formulae, half-base q-calculus, binary telescopes, and discrete extrapolation},
  pdfkeywords={Fabius function, Rvachev up-function, dyadic rational, q-Pochhammer, Gaussian binomial, Thue-Morse, Toeplitz, spline}
}
\pagestyle{fancy}
\fancyhf{}
\fancyhead[L]{\small Fabius--Rvachev synthesis}
\fancyhead[R]{\small\nouppercase{\leftmark}}
\fancyfoot[C]{\thepage}
\renewcommand{\headrulewidth}{0.4pt}
\titleformat{\section}{\Large\bfseries}{\thesection}{0.7em}{}
\titleformat{\subsection}{\large\bfseries}{\thesubsection}{0.7em}{}
\titleformat{\subsubsection}{\normalsize\bfseries}{\thesubsubsection}{0.7em}{}
\setcounter{secnumdepth}{3}
\setcounter{tocdepth}{2}
\setlist[itemize]{leftmargin=2em,itemsep=0.25em,topsep=0.4em}
\setlist[enumerate]{leftmargin=2.3em,itemsep=0.25em,topsep=0.4em}
\setlength{\emergencystretch}{2em}
\newtheorem{theorem}{Theorem}[section]
\newtheorem{proposition}[theorem]{Proposition}
\newtheorem{lemma}[theorem]{Lemma}
\newtheorem{corollary}[theorem]{Corollary}
\newtheorem{conjecture}[theorem]{Conjecture}
\theoremstyle{definition}
\newtheorem{definition}[theorem]{Definition}
\newtheorem{algorithm}[theorem]{Algorithm}
\newtheorem{example}[theorem]{Example}
\theoremstyle{remark}
\newtheorem{remark}[theorem]{Remark}
\newtheorem{warning}[theorem]{Warning}
\newcommand{\R}{\mathbb R}
\newcommand{\C}{\mathbb C}
\newcommand{\Q}{\mathbb Q}
\newcommand{\N}{\mathbb N}
\newcommand{\Z}{\mathbb Z}
\newcommand{\Up}{\operatorname{up}}
\newcommand{\sinc}{\operatorname{sinc}}
\newcommand{\wt}{\operatorname{wt}_2}
\newcommand{\e}{\mathrm e}
\newcommand{\ii}{\mathrm i}
\newcommand{\dd}{\,\mathrm d}
\newcommand{\E}{\mathbb E}
\newcommand{\Prob}{\mathbb P}
\newcommand{\Var}{\operatorname{Var}}
\newcommand{\bigO}{\mathcal O}
\newcommand{\smallo}{o}
\newcommand{\qbinom}[3]{\genfrac{[}{]}{0pt}{}{#1}{#2}_{#3}}
\newcommand{\repo}{\nolinkurl{Analysis/FabiusFunction}}
\newcommand{\coef}[2]{[X^{#1}]\,#2}
\newcommand{\tm}{\varepsilon}
\newcommand{\Dq}{\mathfrak D}
\newcommand{\shiftop}{\mathsf L}
\newenvironment{proofidea}{\par\smallskip\noindent\textit{Proof architecture.}\ }{\hfill$\square$\par\smallskip}
\newenvironment{boxedremark}{%
  \par\smallskip\noindent\begin{center}\begin{minipage}{0.94\textwidth}%
  \color{black}\hrule\vspace{0.6em}\small
}{%
  \vspace{0.6em}\hrule\end{minipage}\end{center}\smallskip
}
\lstdefinestyle{pseudo}{
  basicstyle=\ttfamily\small,
  backgroundcolor=\color{shadegray},
  frame=single,
  rulecolor=\color{rulegray},
  columns=fullflexible,
  keepspaces=true,
  showstringspaces=false,
  xleftmargin=0.7em,
  xrightmargin=0.7em,
  aboveskip=0.8em,
  belowskip=0.8em
}

\title{\bfseries The Dyadic Web of the Fabius Function\\[0.35em]
\large Moment corrections, half-base $q$-interpolation, binary telescopes,\\
and discrete extrapolation}
\author{A complementary exposition to \emph{The Fabius Function and Rvachev's Up-Function}\\
\small based on Vladimir Reshetnikov's \texttt{ProveIt} formal development}
\date{Repository snapshot: 25 August 2026}

\begin{document}
\maketitle

\begin{abstract}
Exact values of the Fabius function at dyadic rationals appear in several forms that look
unrelated: a Thue--Morse sum corrected by the even moments of Rvachev's up-function, a finite
half-base $q$-binomial expression with a seemingly arbitrary complex translation, a
bit-recursive Taylor reduction, an absolutely convergent binary series at arbitrary argument,
and a finite $q$-weighted discrete approximant whose limit is the signed global Fabius
function.  This companion article identifies the common algebra behind these formulae and
makes their connections explicit.

The central object is the coefficient
\[
 [X^n]\bigl(B_m(X)\mathscr A(X)\bigr),
 \qquad
 B_m(X)=\sum_{h<m}(-1)^{\wt(h)}\e^{(m-1-h)X},
\]
where $\mathscr A$ is the exponential moment series of the Fabius probability law.  Direct
expansion gives the classical moment formula; probabilistic interpretation gives an
expectation formula; Cauchy's formula gives a contour representation; and geometric divided
differences through $-1,-2,\ldots,-2^n$ give the $q$-binomial formula.  The same coefficient
also produces new finite identities: an exact moment-corrected expansion in centered splines,
its reciprocal deconvolution, and a residue formula for the Gaussian interpolation row.

At arbitrary arguments the two infinite constructions have different architectures.  The
binary series is a genuine telescoping sum whose summands are individually independent of the
complex shift.  The discrete formula is instead a Toeplitz transform of shifted splines; its
finite rows generally depend on the shift, although their limit does not.  Writing the row
generating polynomial as
\[
 W_n(z)=\frac{(z/2;1/2)_n}{(1/2;1/2)_n}
\]
reveals it as a $q$-Richardson filter: it preserves constants and annihilates the geometric
error modes $2^{-p},2^{-2p},\ldots,2^{-np}$.  This yields an explicit uniform
$\bigO(4^{-n})$ error bound.  Most importantly, at a dyadic input $m/2^s$ the discrete
approximant stabilizes exactly for every row $n\ge s$; thus the ``limit formula'' literally
becomes the exact dyadic $q$-formula after the denominator has been resolved.  A synchronized
nearest-dyadic version provides a second, termwise translation-invariant discrete limit.
\end{abstract}

\tableofcontents
\newpage

\section{Scope, conventions, and the formula map}
\label{sec:scope}

This article is a companion rather than a replacement for the comprehensive exposition
\cite{mainarticle}.  It assumes the construction and basic analytic theory developed there
and concentrates on one cluster of results: exact dyadic evaluation, binary reduction, finite
half-base $q$-calculus, and discrete limits.  The proofs formalized in Lean under \repo
\cite{proveit} are the source of the exact identities used below.  Several consequences in
this article are then derived directly from those identities; a formalization map is given in
\cref{sec:formal-map}.

We use three related functions, with the same convention as the companion article:
\begin{itemize}
\item $F$ is the bounded Fabius distribution function, equal to $0$ on $(-\infty,0]$ and to
      $1$ on $[1,\infty)$;
\item $\Up$ is the even compactly supported Rvachev density on $[-1,1]$;
\item $\mathcal F$ is the signed global extension satisfying
      $\mathcal F'(x)=2\mathcal F(2x)$ on $\R$.
\end{itemize}
They agree on $[0,1]$.  Put
\[
 \tm_r=(-1)^{\wt(r)},\qquad
 c_k=\int_{-1}^{1}u^{2k}\Up(u)\dd u,
\]
and let $d_n$ denote the half moments normalized by
\begin{equation}
 F(2^{-n})=2^{-\binom n2}\frac{d_n}{n!}.                 \label{eq:inverse-dyadic-input}
\end{equation}
The normalized moment series is
\begin{equation}
 \mathscr A(X)=\sum_{n\ge0}d_n\frac{X^n}{n!}.             \label{eq:A-def}
\end{equation}
We reserve $q$ for the fixed base $q=1/2$ and use $Q\in\C$ for a freely chosen translation.
This distinction is essential: in the expressions $(a;q)_n$ and
$(r-2^k+Q)^{n+k}$ the two symbols play completely different roles.

\subsection{Five presentations of the same value}

The identities studied below fall into five families.  Their superficial differences are
summarized in \cref{tab:formula-map}; the rest of the article explains the arrows between
rows.

\begin{table}[ht]
\centering
\small
\begin{tabularx}{\textwidth}{@{}>{\raggedright\arraybackslash}p{0.19\textwidth}
>{\raggedright\arraybackslash}p{0.18\textwidth}
>{\raggedright\arraybackslash}p{0.22\textwidth}X@{}}
\toprule
Presentation & Range and status & Principal coefficients & Mechanism \\
\midrule
Moment/Thue--Morse formula & finite and exact at $m/2^n$ & $c_k$ and $\tm_h$ & coefficient expansion in the centered probability law \\
Half-base $q$-formula & finite and exact at $m/2^n$ & $\qbinom nk{1/2}$ and $(1/2;1/2)_n$ & geometric divided difference after block cancellation \\
Binary Taylor series & infinite and absolutely convergent for $x\ge0$; finite at dyadics & binary prefixes and $d_j$ & exact residual telescope along the binary expansion \\
Global $q$-binary series & same outer series as the preceding row & a finite $q$-identity inside each scale & termwise replacement of a Taylor polynomial \\
Discrete $q$-limit & one finite row for each $n$, convergent for all $x\in\R$ & normalized Gaussian row $\omega_{n,j}$ & Toeplitz transform of shifted finite splines \\
\bottomrule
\end{tabularx}
\caption{The formula families.  The two infinite constructions are not partial-sum versions
of one another.}
\label{tab:formula-map}
\end{table}

The main structural conclusions can be stated before any calculation.

\begin{boxedremark}
\textbf{The common core.}
The exact dyadic formulae are all coefficient extractors applied to the same entire function
$B_m\mathscr A$.  The binary series repeatedly reuses its Taylor coefficients at successive
binary scales.  The discrete formula applies a normalized $q$-Pochhammer polynomial to the
sequence of finite-spline degrees.  At a dyadic input these two arbitrary-argument mechanisms
both become finite: the binary series terminates, while the discrete rows stabilize exactly.
\end{boxedremark}

\subsection{The exact dyadic formula taken as input}

For $0\le m\le2^n$ the moment formula proved in the formal development is
\begin{equation}
 \boxed{
 F\!\left(\frac{m}{2^n}\right)
 =\frac{2^{-\binom{n+1}{2}}}{n!}
 \sum_{h=0}^{m-1}\tm_h
 \sum_{k=0}^{\lfloor n/2\rfloor}
 \binom n{2k}c_k(2m-2h-1)^{n-2k}.}                        \label{eq:input-moment-dyadic}
\end{equation}
The same finite expression, interpreted with $\mathcal F$ on the left, is valid for every
$m,n\in\N$.  The $q$-binomial version is
\begin{align}
 \mathcal F(m/2^n)
 &=\frac{1}{2^{n^2}(1/2;1/2)_n}
 \sum_{k=0}^{n}
 \frac{\qbinom nk{1/2}}{4^{\binom k2}(n+k)!}             \notag\\
 &\qquad\times
 \sum_{r=0}^{m2^k-1}\tm_r(r-m2^k+Q)^{n+k},               \label{eq:input-q-dyadic}
\end{align}
valid for every $Q\in\C$ and independent of $Q$ as a complete finite expression.  A central
goal below is to explain why \eqref{eq:input-moment-dyadic} and
\eqref{eq:input-q-dyadic} are two readings of one coefficient.

\section{The master coefficient and its probabilistic meaning}
\label{sec:master}

\subsection{The probability law behind the moment series}

Let $U_1,U_2,\ldots$ be independent random variables, each uniformly distributed on
$[0,1]$, and set
\begin{equation}
 Y=\sum_{j=1}^{\infty}\frac{U_j}{2^j},
 \qquad Z=2Y-1.                                             \label{eq:Y-Z-def}
\end{equation}
Then $Y\in[0,1]$, its distribution function is $F$, and $Z$ has density $\Up$ on $[-1,1]$.
Consequently
\begin{equation}
 d_n=\E(Y^n),\qquad c_k=\E(Z^{2k}),\qquad \E(Z^{2k+1})=0. \label{eq:prob-moments}
\end{equation}
The moment generating function of one uniform variable is
\[
 E_0(X)=\frac{\e^X-1}{X}=\int_0^1\e^{Xu}\dd u,
\]
with the removable value $E_0(0)=1$.  Independence gives the locally uniformly convergent
entire product
\begin{equation}
 \mathscr A(X)=\E(\e^{XY})
 =\prod_{j=1}^{\infty}E_0(X/2^j).                          \label{eq:A-product}
\end{equation}
Separating the first factor, or merely replacing $X$ by $2X$ in the product, gives the
refinement equation
\begin{equation}
 \mathscr A(2X)=E_0(X)\mathscr A(X).                       \label{eq:A-refinement}
\end{equation}
The centered moment series
\begin{equation}
 \mathscr C(T)=\E(\e^{TZ})
 =\sum_{k\ge0}c_k\frac{T^{2k}}{(2k)!}                     \label{eq:C-def}
\end{equation}
satisfies the exact change of coordinates
\begin{equation}
 \boxed{\mathscr A(X)=\e^{X/2}\mathscr C(X/2).}           \label{eq:A-C-transform}
\end{equation}
The familiar hyperbolic-sine form follows by writing
$E_0(t)=\e^{t/2}\sinh(t/2)/(t/2)$:
\begin{align}
 \mathscr A(X)
 &=\e^{X/2}\prod_{j=1}^{\infty}
   \frac{\sinh(X/2^{j+1})}{X/2^{j+1}},                    \label{eq:A-sinh-product}\\
 \mathscr C(T)
 &=\prod_{j=1}^{\infty}\frac{\sinh(T/2^j)}{T/2^j}.        \label{eq:C-sinh-product}
\end{align}
Thus the moment series used in exact arithmetic is the real-exponential counterpart of the
sinc product in Fourier analysis.

\subsection{The Thue--Morse prefix series}

For $m\in\N$ define
\begin{equation}
 B_m(X)=\sum_{h=0}^{m-1}\tm_h\e^{(m-1-h)X},
 \qquad B_0(X)=0.                                         \label{eq:Bm-def}
\end{equation}
This finite exponential polynomial records the numerator $m$ and nothing else.  Its two
factors have distinct roles:
\begin{itemize}
\item $B_m$ carries the finite Thue--Morse prefix and therefore the binary numerator;
\item $\mathscr A$ carries the universal Fabius moments and therefore the analytic law.
\end{itemize}
Their product is the promised master object.

\begin{theorem}[Master coefficient identity]
\label{thm:master-coefficient}
For every $m,n\in\N$,
\begin{equation}
 \boxed{
 \mathcal F\!\left(\frac{m}{2^n}\right)
 =2^{-\binom n2}[X^n]\bigl(B_m(X)\mathscr A(X)\bigr).}    \label{eq:master-coefficient}
\end{equation}
If $0\le m\le2^n$, the left side is the bounded value $F(m/2^n)$.
\end{theorem}

\begin{proof}
Put $t_h=2m-2h-1$.  By \eqref{eq:C-def},
\[
 \sum_{k=0}^{\lfloor n/2\rfloor}\binom n{2k}c_kt_h^{n-2k}
 =n![W^n]\bigl(\e^{t_hW}\mathscr C(W)\bigr).
\]
Using \eqref{eq:A-C-transform} at $X=2W$ gives
$\mathscr C(W)=\e^{-W}\mathscr A(2W)$.  Since $t_h-1=2(m-1-h)$,
\begin{align*}
 n![W^n]\bigl(\e^{t_hW}\mathscr C(W)\bigr)
 &=n![W^n]\bigl(\e^{2(m-1-h)W}\mathscr A(2W)\bigr)\\
 &=n!2^n[X^n]\bigl(\e^{(m-1-h)X}\mathscr A(X)\bigr).
\end{align*}
Multiply by $\tm_h$, sum over $h<m$, and insert the result in
\eqref{eq:input-moment-dyadic}.  The power of two becomes
$2^{n-\binom{n+1}{2}}=2^{-\binom n2}$.
\end{proof}

The theorem compresses the double sum in \eqref{eq:input-moment-dyadic} into one coefficient,
but it also expands naturally in several different directions.

\begin{corollary}[Expectation form]
\label{cor:expectation-dyadic}
For every $m,n\in\N$,
\begin{align}
 \mathcal F(m/2^n)
 &=\frac{2^{-\binom n2}}{n!}
   \sum_{h=0}^{m-1}\tm_h\,
   \E\bigl(m-1-h+Y\bigr)^n                              \label{eq:expectation-Y}\\
 &=\frac{2^{-\binom{n+1}{2}}}{n!}
   \sum_{h=0}^{m-1}\tm_h\,
   \E\bigl(2m-2h-1+Z\bigr)^n.                            \label{eq:expectation-Z}
\end{align}
\end{corollary}

\begin{proof}
Because $\mathscr A(X)=\E(\e^{XY})$,
\[
 [X^n]\bigl(\e^{aX}\mathscr A(X)\bigr)
 =\frac1{n!}\E(a+Y)^n.
\]
Apply this termwise to $B_m$.  The second form follows from
$m-1-h+Y=(2m-2h-1+Z)/2$.
\end{proof}

Expanding the centered expectation in \eqref{eq:expectation-Z}, using symmetry to discard the
odd moments, recovers \eqref{eq:input-moment-dyadic} immediately.  In this sense the
Thue--Morse/moment formula is simply the binomial theorem applied inside one probability
expectation.

\begin{corollary}[Cauchy coefficient representation]
\label{cor:cauchy-X}
For any $\rho>0$,
\begin{equation}
 \mathcal F(m/2^n)
 =\frac{2^{-\binom n2}}{2\pi\ii}
  \int_{|X|=\rho}\frac{B_m(X)\mathscr A(X)}{X^{n+1}}\dd X. \label{eq:cauchy-X}
\end{equation}
\end{corollary}

This contour formula is elementary, but useful conceptually: the $q$-binomial formula will
turn out to be a second residue computation, now in an auxiliary geometric node variable.

\subsection{The numerator and the law separate completely}

Expanding \eqref{eq:master-coefficient} without centering gives another exact finite form:
\begin{equation}
 \mathcal F(m/2^n)
 =\frac{2^{-\binom n2}}{n!}
 \sum_{h=0}^{m-1}\tm_h
 \sum_{j=0}^{n}\binom njd_j(m-1-h)^{n-j}.                 \label{eq:d-moment-prefix}
\end{equation}
The $d_j$ are universal and may be precomputed once.  The numerator enters only through the
finite power sums
\(
 \sum_{h<m}\tm_h(m-1-h)^r.
\)
This is often a more transparent bridge to the Taylor reduction series than the centered
$c_k$ formula, because the reduction polynomials themselves are written in the $d_j$.

\begin{remark}[Two independent cancellations]
The Thue--Morse prefix in \eqref{eq:d-moment-prefix} may annihilate low ordinary powers when
$m$ is a complete binary block.  The $q$-binomial formula introduces a second cancellation
across the geometric scale index.  These cancellations live in different variables and should
not be conflated: one acts on $h$ or $r$ inside a binary block, the other on the scale node
$-2^k$.
\end{remark}

\section{Moment coordinates and alternative exact representations}
\label{sec:moment-alternatives}

\subsection[The triangular transforms between the c and d sequences]{The triangular transforms between $c_k$ and $d_n$}

The affine relation $Z=2Y-1$ gives both directions of the moment change of coordinates.

\begin{proposition}[Moment transforms]
\label{prop:moment-transforms}
For every $n,k\in\N$,
\begin{align}
 d_n
 &=2^{-n}\sum_{k=0}^{\lfloor n/2\rfloor}\binom n{2k}c_k,
                                                               \label{eq:d-from-c-new}\\
 c_k
 &=\sum_{j=0}^{2k}(-1)^j2^j\binom{2k}{j}d_j.              \label{eq:c-from-d-new}
\end{align}
\end{proposition}

\begin{proof}
Since $Y=(1+Z)/2$, expand $Y^n$ and use the vanishing odd moments of $Z$.  Conversely expand
$(2Y-1)^{2k}$; because $2k$ is even, the sign of the term of degree $j$ is $(-1)^j$.
\end{proof}

The first formula explains the odd-looking shift $2m-2h-1$ in
\eqref{eq:input-moment-dyadic}: it is precisely the centered coordinate obtained after
converting the half moments to the even moments of $\Up$.

\subsection{An infinite-cube composition formula}

Each factor in \eqref{eq:A-product} has the expansion
\[
 E_0(X/2^j)=\sum_{r\ge0}\frac{X^r}{2^{jr}(r+1)!}.
\]
Taking the coefficient of $X^n$ gives an exact positive series over finite-support weak
compositions.

\begin{proposition}[Composition representation of the half moments]
\label{prop:composition-d}
For every $n\in\N$,
\begin{equation}
 d_n=n!\!
 \sum_{\substack{r_1,r_2,\ldots\ge0\\r_1+r_2+\cdots=n}}
 \prod_{j\ge1}\frac{2^{-jr_j}}{(r_j+1)!}.                 \label{eq:composition-d}
\end{equation}
Only finitely many $r_j$ are nonzero in each summand, and the series converges absolutely.
Equivalently,
\begin{equation}
 d_n=\lim_{J\to\infty}\int_{[0,1]^J}
       \left(\sum_{j=1}^{J}\frac{u_j}{2^j}\right)^n
       \dd u_1\cdots\dd u_J.                              \label{eq:cube-d}
\end{equation}
\end{proposition}

This gives a $q$-free, manifestly positive representation of every inverse-dyadic value:
\begin{equation}
 F(2^{-n})=\frac{2^{-\binom n2}}{n!}d_n
 =2^{-\binom n2}
 \sum_{r_1+r_2+\cdots=n}
 \prod_{j\ge1}\frac{2^{-jr_j}}{(r_j+1)!}.                \label{eq:inverse-composition}
\end{equation}

\subsection{Cumulants, Bernoulli numbers, and Bell polynomials}

The infinite product also gives a compact partition formula.  With Bernoulli numbers in the
convention $B_2=1/6$, $B_4=-1/30$, one has
\begin{equation}
 \log\frac{\e^X-1}{X}
 =\frac X2+\sum_{r\ge1}\frac{B_{2r}}{2r(2r)!}X^{2r}.      \label{eq:log-exprel}
\end{equation}
Summing this identity over the dyadic scales in \eqref{eq:A-product} yields
\begin{equation}
 \log\mathscr A(X)
 =\frac X2+
  \sum_{r\ge1}
  \frac{B_{2r}}{2r(2r)!(2^{2r}-1)}X^{2r}.                 \label{eq:log-A}
\end{equation}
Thus the cumulants $\kappa_j$ of $Y$ are
\begin{equation}
 \kappa_1=\frac12,\qquad
 \kappa_{2r}=\frac{B_{2r}}{2r(2^{2r}-1)},\qquad
 \kappa_{2r+1}=0\quad(r\ge1).                            \label{eq:Y-cumulants}
\end{equation}

\begin{theorem}[Bell-polynomial representation]
\label{thm:bell-d}
Let $\mathcal B_n$ be the complete exponential Bell polynomial.  Then
\begin{equation}
 d_n=\mathcal B_n(\kappa_1,\ldots,\kappa_n),               \label{eq:bell-d}
\end{equation}
and explicitly
\begin{equation}
 d_n=n!\!
 \sum_{\substack{m_1,\ldots,m_n\ge0\\
                  1m_1+2m_2+\cdots+nm_n=n}}
 \prod_{r=1}^{n}\frac1{m_r!}
 \left(\frac{\kappa_r}{r!}\right)^{m_r}.                 \label{eq:partition-d}
\end{equation}
Consequently
\begin{equation}
 \boxed{
 F(2^{-n})=\frac{2^{-\binom n2}}{n!}
 \mathcal B_n(\kappa_1,\ldots,\kappa_n).}                 \label{eq:inverse-bell}
\end{equation}
\end{theorem}

\begin{proof}
Equation \eqref{eq:log-A} is the standard cumulant expansion
$\log\mathscr A(X)=\sum_{j\ge1}\kappa_jX^j/j!$.  Exponentiating and comparing coefficients
is exactly the defining generating identity for the complete Bell polynomials.
\end{proof}

The representation is useful when Bernoulli arithmetic is already available.  It also gives
the moment recurrence
\begin{equation}
 d_{n+1}=\sum_{j=0}^{n}\binom nj\kappa_{j+1}d_{n-j},       \label{eq:cumulant-moment-recurrence}
\end{equation}
while \eqref{eq:A-refinement} gives the rational recurrence
\begin{equation}
 (2^n-1)d_n
 =\sum_{r=1}^{n}\binom nr\frac{d_{n-r}}{r+1}
 \qquad(n\ge1).                                           \label{eq:d-recurrence-companion}
\end{equation}

\begin{table}[ht]
\centering
\small
\begin{tabular}{@{}c c c c@{}}
\toprule
$n$ & $d_n$ & $c_n$ & $F(2^{-n})$ \\
\midrule
0 & $1$ & $1$ & $1$ \\
1 & $1/2$ & $1/9$ & $1/2$ \\
2 & $5/18$ & $19/675$ & $5/72$ \\
3 & $1/6$ & $583/59535$ & $1/288$ \\
4 & $143/1350$ & $132809/32531625$ & $143/2073600$ \\
5 & $19/270$ & $46840699/24405225075$ & $19/33177600$ \\
\bottomrule
\end{tabular}
\caption{The columns use different index conventions: $c_n$ is the $2n$th centered moment,
whereas $d_n$ is the $n$th half moment.}
\label{tab:moments}
\end{table}

\subsection{Exact moment correction of the matched spline}
\label{subsec:moment-corrected-spline}

The centered finite spline of degree $p\ge1$ is
\begin{equation}
 S_p(x)=\frac{(-1)^p}{2^{\binom p2}p!}
 \sum_{r=0}^{N_p(x)-1}\tm_r\left(r-2^px+\frac12\right)^p,
 \qquad
 N_p(x)=\max\!\left(0,\left\lfloor2^px+\frac12\right\rfloor\right). \label{eq:centered-spline-local}
\end{equation}
At degree zero set
\begin{equation}
 S_0(x)=\sum_{r=0}^{N_0(x)-1}\tm_r.                        \label{eq:S0-local}
\end{equation}
The companion article proves the global estimate
\begin{equation}
 \sup_{x\in\R}|S_p(x)-\mathcal F(x)|\le2^{-p}.             \label{eq:centered-global-error-input}
\end{equation}
At the matched grid point $m/2^p$ the cutoff is $m$ and
\begin{equation}
 S_p(m/2^p)=\frac{1}{2^{\binom{p+1}{2}}p!}
 \sum_{h=0}^{m-1}\tm_h(2m-2h-1)^p.                       \label{eq:matched-spline}
\end{equation}
The $k=0$ part of \eqref{eq:input-moment-dyadic} is exactly this matched spline.  The remaining
moments supply a finite correction hierarchy.

\begin{theorem}[Moment-corrected matched-spline identity]
\label{thm:moment-corrected-spline}
For all $m,n\in\N$,
\begin{equation}
 \boxed{
 \mathcal F(m/2^n)=
 \sum_{k=0}^{\lfloor n/2\rfloor}
 \frac{c_k}{(2k)!\,2^{k(2n-2k+1)}}
 S_{n-2k}\!\left(\frac{m}{2^{n-2k}}\right).}              \label{eq:moment-corrected-spline}
\end{equation}
For $0\le m\le2^n$ the left side is $F(m/2^n)$.  The spline arguments on the right are
understood globally and may lie outside $[0,1]$.
\end{theorem}

\begin{proof}
In \eqref{eq:input-moment-dyadic} put $p=n-2k$.  By
\eqref{eq:matched-spline},
\[
 \sum_{h=0}^{m-1}\tm_h(2m-2h-1)^p
 =2^{\binom{p+1}{2}}p!\,S_p(m/2^p).
\]
Moreover
\[
 \frac1{n!}\binom n{2k}p!=\frac1{(2k)!},
\]
and
\[
 \binom{n+1}{2}-\binom{p+1}{2}=k(2n-2k+1).
\]
Substitution gives \eqref{eq:moment-corrected-spline}.
\end{proof}

The formula has three noteworthy features.
\begin{enumerate}
\item Its first term is the obvious finite-spline approximation $S_n(m/2^n)$.
\item The correction drops the spline degree by two at each step, reflecting the evenness of
      $\Up$ and the disappearance of odd centered moments.
\item The numerator $m$ remains fixed, while the denominator is coarsened from $2^n$ to
      $2^{n-2k}$.  Thus the formula couples degree reduction to dyadic rescaling.
\end{enumerate}

\begin{corollary}[Exact defect of the matched spline]
\label{cor:matched-defect}
For every $m,n\in\N$,
\begin{equation}
 \mathcal F(m/2^n)-S_n(m/2^n)
 =\sum_{k=1}^{\lfloor n/2\rfloor}
 \frac{c_k}{(2k)!\,2^{k(2n-2k+1)}}
 S_{n-2k}\!\left(\frac{m}{2^{n-2k}}\right).               \label{eq:matched-defect}
\end{equation}
\end{corollary}

This is stronger than an error estimate: it identifies the error exactly as a finite sequence
of lower-degree global splines.  The small factor in the first correction is
$c_1/(2\cdot2^{2n-1})=1/(9\cdot2^{2n})$.

\subsection{The reciprocal correction transform}

The previous relation is triangular and can be inverted.  Define reciprocal centered moments
$\gamma_k$ by
\begin{equation}
 \frac1{\mathscr C(T)}=
 \sum_{k\ge0}\gamma_k\frac{T^{2k}}{(2k)!}.                \label{eq:gamma-def}
\end{equation}
Thus
\begin{equation}
 \gamma_0=1,
 \qquad
 \gamma_n=-\sum_{k=0}^{n-1}\binom{2n}{2k}\gamma_kc_{n-k}
 \quad(n\ge1),                                             \label{eq:gamma-recurrence}
\end{equation}
and the first values are
\begin{equation}
 1,\quad-\frac19,\quad\frac{31}{675},\quad
 -\frac{2347}{59535},\quad\frac{1904369}{32531625},\ldots. \label{eq:gamma-values}
\end{equation}

\begin{theorem}[Reciprocal moment deconvolution]
\label{thm:reciprocal-spline}
For every $m,n\in\N$,
\begin{equation}
 \boxed{
 S_n(m/2^n)=
 \sum_{k=0}^{\lfloor n/2\rfloor}
 \frac{\gamma_k}{(2k)!\,2^{k(2n-2k+1)}}
 \mathcal F\!\left(\frac{m}{2^{n-2k}}\right).}            \label{eq:reciprocal-spline}
\end{equation}
\end{theorem}

\begin{proof}
For a fixed numerator $m$, introduce
\[
 P_p(m)=\sum_{h=0}^{m-1}\tm_h(2m-2h-1)^p,
 \qquad
 G_p(m)=2^{\binom{p+1}{2}}p!\,\mathcal F(m/2^p).
\]
The exact moment formula says, in exponential generating functions,
\[
 \sum_{p\ge0}G_p(m)\frac{T^p}{p!}
 =\mathscr C(T)\sum_{p\ge0}P_p(m)\frac{T^p}{p!}.
\]
Multiply by $1/\mathscr C(T)$, compare coefficients, and then use
$P_p(m)=2^{\binom{p+1}{2}}p!S_p(m/2^p)$.  The same triangular exponent simplification as in
\cref{thm:moment-corrected-spline} gives \eqref{eq:reciprocal-spline}.
\end{proof}

Equations \eqref{eq:moment-corrected-spline} and \eqref{eq:reciprocal-spline} are a finite
convolution/deconvolution pair.  They show that the exact dyadic values and the matched centered
splines carry precisely the same information once the universal moment series $\mathscr C$ is
known.

\begin{boxedremark}
\textbf{First bridge to the discrete theory.}
The centered spline is not merely an approximation from which exact dyadic arithmetic is
separate.  At a matched grid point, the exact value is obtained by a finite universal moment
correction of that spline, and the correction is invertible.  Later, the half-base
$q$-Pochhammer row will give a second exact correction, using consecutive spline degrees at
the \emph{same} argument rather than parity-separated degrees at rescaled arguments.
\end{boxedremark}

\section[The half-base q-formula as geometric coefficient extraction]{The half-base $q$-formula as geometric coefficient extraction}
\label{sec:q-extraction}

\subsection{The Gaussian row is a divided difference}

For a general base $q$, write
\begin{equation}
 (a;q)_n=\prod_{j=0}^{n-1}(1-aq^j),                       \label{eq:qpoch-def-local}
\end{equation}
and let $\qbinom nkq$ denote the Gaussian binomial coefficient.  At $q=1/2$ the finite
$q$-binomial theorem is
\begin{equation}
 (z;1/2)_n=\sum_{k=0}^{n}(-1)^k2^{-\binom k2}
 \qbinom nk{1/2}z^k.                                      \label{eq:finite-q-binomial-local}
\end{equation}
Put
\begin{equation}
 x_k=-2^k,
 \qquad
 w_{n,k}=(-1)^k2^{-\binom k2}\qbinom nk{1/2},
 \qquad
 \mathfrak m_n=\prod_{\ell=1}^{n}(2^\ell-1).              \label{eq:q-nodes-weights}
\end{equation}
Then
\begin{equation}
 \sum_{k=0}^{n}w_{n,k}x_k^d=
 \begin{cases}
 0,&0\le d<n,\\
 \mathfrak m_n,&d=n,
 \end{cases}                                               \label{eq:q-node-moments}
\end{equation}
and
\begin{equation}
 \frac{w_{n,k}}{\mathfrak m_n}
 =\frac1{\prod_{\ell\ne k}(x_k-x_\ell)},
 \qquad
 (1/2;1/2)_n=2^{-\binom{n+1}{2}}\mathfrak m_n.             \label{eq:q-barycentric-normalization}
\end{equation}
Thus the functional
\begin{equation}
 \Dq_n[P]=\frac1{\mathfrak m_n}\sum_{k=0}^{n}w_{n,k}P(-2^k) \label{eq:Dq-def}
\end{equation}
is the divided difference $[-1,-2,\ldots,-2^n]P$.  For every polynomial of degree at most
$n$, it returns the leading coefficient.

\begin{proposition}[Residue form of the geometric extractor]
\label{prop:q-residue-extractor}
Let
\[
 \Pi_n(z)=\prod_{k=0}^{n}(z+2^k),
\]
and let $P$ be a polynomial of degree at most $n$.  If $\Gamma$ encloses all nodes
$-1,-2,\ldots,-2^n$, then
\begin{equation}
 [z^n]P(z)=\Dq_n[P]
 =\frac{1}{2\pi\ii}\int_\Gamma\frac{P(z)}{\Pi_n(z)}\dd z. \label{eq:q-residue-extractor}
\end{equation}
\end{proposition}

\begin{proof}
The residues at the simple poles $x_k=-2^k$ are
$P(x_k)/\Pi_n'(x_k)=P(x_k)/\prod_{\ell\ne k}(x_k-x_\ell)$, giving the barycentric sum in
\eqref{eq:Dq-def}.  Lagrange interpolation identifies that sum with the leading coefficient.
\end{proof}

The $q$-symbols in the dyadic formula therefore have a concrete and limited task: they are the
closed form of a divided-difference stencil on a geometric grid.  No infinite
basic-hypergeometric summation is involved.

\subsection{The parametric coefficient polynomial}

Introduce an auxiliary node variable $z$ and define
\begin{equation}
 \mathscr R_m(z;X)
 =\frac{\e^{-X}B_m(zX)\mathscr A(zX)}{\mathscr A(-X)},
 \qquad
 r_{m,n}(z)=[X^n]\mathscr R_m(z;X).                        \label{eq:Rmn-def}
\end{equation}
Because $\mathscr A(-X)$ is a unit with constant term $1$, this is a well-defined formal
series and an entire analytic germ.  More importantly, $r_{m,n}$ is a polynomial in $z$.

\begin{lemma}[Node polynomial and its leading coefficient]
\label{lem:Rmn-leading}
For fixed $m,n$, the polynomial $r_{m,n}(z)$ has degree at most $n$ and
\begin{equation}
 [z^n]r_{m,n}(z)=[X^n]\bigl(B_m(X)\mathscr A(X)\bigr)
 =2^{\binom n2}\mathcal F(m/2^n).                         \label{eq:Rmn-leading}
\end{equation}
\end{lemma}

\begin{proof}
Write
\(
 \e^{-X}/\mathscr A(-X)=\sum_{j\ge0}u_jX^j
\)
with $u_0=1$.  Since
$B_m(zX)\mathscr A(zX)=(B_m\mathscr A)(zX)$, the coefficient of $X^n$ is a polynomial in $z$
of degree at most $n$.  Its $z^n$ term can only use the constant $u_0$, so its leading
coefficient is $[X^n](B_m\mathscr A)$.  Apply \cref{thm:master-coefficient}.
\end{proof}

The exact dyadic value is now the leading coefficient of a degree-$n$ polynomial sampled at
$-1,-2,\ldots,-2^n$.  Applying \cref{prop:q-residue-extractor} gives a compact alternative
to the displayed $q$-sum.

\begin{corollary}[Geometric-node contour representation]
\label{cor:q-node-contour}
For every $m,n\in\N$,
\begin{equation}
 \boxed{
 \mathcal F(m/2^n)=
 \frac{2^{-\binom n2}}{2\pi\ii}
 \int_\Gamma
 \frac{r_{m,n}(z)}{\prod_{k=0}^{n}(z+2^k)}\dd z,}          \label{eq:q-node-contour}
\end{equation}
where $\Gamma$ encloses the geometric nodes.
\end{corollary}

Together, \eqref{eq:cauchy-X} and \eqref{eq:q-node-contour} exhibit two independent Cauchy
extractions:
\begin{center}
\begin{tabular}{c@{\qquad}c}
coefficient in the analytic variable $X$
& leading coefficient in the geometric node $z$ \\
$[X^n](B_m\mathscr A)$
& $[z^n]r_{m,n}(z)$ \\
ordinary Cauchy kernel $X^{-n-1}$
& barycentric kernel $\Pi_n(z)^{-1}$.
\end{tabular}
\end{center}
The finite $q$-binomial formula is what results when the second contour is evaluated by
residues and each node value is converted to a Thue--Morse block sum.

\subsection{Why the node values are Thue--Morse power sums}

For a complete binary block set
\begin{equation}
 C_k(X)=\sum_{r=0}^{2^k-1}\tm_r\e^{(r-2^k)X}.              \label{eq:Ck-local}
\end{equation}
The digit factorization
\(
 \sum_{r<2^k}\tm_rz^r=\prod_{j=0}^{k-1}(1-z^{2^j})
\)
gives
\begin{equation}
 C_k(X)=(-1)^k2^{\binom k2}X^k
 \e^{-X}\prod_{j=0}^{k-1}E_0(-2^jX).                     \label{eq:Ck-factor-local}
\end{equation}
Hence $C_k$ has a zero of exact order $k$.  The refinement
\eqref{eq:A-refinement}, iterated $k$ times, gives
\begin{equation}
 \frac{\e^{-X}\mathscr A(-2^kX)}{\mathscr A(-X)}
 =\e^{-X}\prod_{j=0}^{k-1}E_0(-2^jX).                    \label{eq:A-geometric-bridge}
\end{equation}
Combining \eqref{eq:Ck-factor-local} and \eqref{eq:A-geometric-bridge} identifies the node
specialization of \eqref{eq:Rmn-def} with a normalized complete block.

For a general numerator, the long block factors as
\begin{equation}
 \sum_{r=0}^{m2^k-1}\tm_r\e^{(r-m2^k)X}
 =B_m(-2^kX)C_k(X).                                       \label{eq:long-block-local}
\end{equation}
After division by the forced factor
$(-1)^k2^{\binom k2}X^k$, the coefficient of $X^n$ is
\begin{equation}
 \frac{(-1)^k2^{-\binom k2}}{(n+k)!}
 \sum_{r=0}^{m2^k-1}\tm_r(r-m2^k)^{n+k}.                 \label{eq:node-value-power-sum}
\end{equation}
Multiplication by the interpolation weight $w_{n,k}$ cancels the sign and contributes a
second factor $2^{-\binom k2}$.  Their product is the characteristic
$4^{-\binom k2}$ of \eqref{eq:input-q-dyadic}.  Finally
\eqref{eq:q-barycentric-normalization} converts $\mathfrak m_n$ into
$2^{n^2}(1/2;1/2)_n$ after the factor $2^{\binom n2}$ from
\cref{lem:Rmn-leading} is included.  This recovers the exact $q$-formula with $Q=0$.

\subsection{Translation invariance is lower-degree annihilation}

The shift $Q$ enters a block generating function by multiplication with $\e^{QX}$.  If
\(
 P_{m,k}(X)
\)
denotes the normalized long-block series after its forced $X^k$ has been removed, then the
$q$-numerator is a weighted sum of $[X^n](\e^{QX}P_{m,k}(X))$.  Expanding the exponential,
\begin{equation}
 [X^n](\e^{QX}P_{m,k}(X))
 =\sum_{s=0}^{n}\frac{Q^{n-s}}{(n-s)!}[X^s]P_{m,k}(X).     \label{eq:Q-expansion}
\end{equation}
For each $s<n$, the Gaussian row sees a polynomial of geometric-node degree at most $s$ and
annihilates it.  Only the $s=n$ term survives; it has no $Q$ factor.

\begin{proposition}[Finite translation invariance]
\label{prop:finite-Q-invariance}
For fixed $m,n$, the polynomial
\begin{equation}
 Q\longmapsto
 \sum_{k=0}^{n}
 \frac{\qbinom nk{1/2}}{4^{\binom k2}(n+k)!}
 \sum_{r=0}^{m2^k-1}\tm_r(r-m2^k+Q)^{n+k}                 \label{eq:Q-polynomial}
\end{equation}
is constant.
\end{proposition}

This is stronger than saying that a limit is independent of $Q$: every coefficient of
$Q,Q^2,\ldots$ vanishes in the finite expression itself.  The distinction becomes important
in \cref{sec:discrete}, where the finite discrete rows generally do depend on $Q$.

\subsection{A useful abstract summary}

The exact $q$-formula is a composition of two annihilators:
\begin{enumerate}
\item the complete Thue--Morse block has a zero of order $k$, removing all inner polynomial
      degrees below $k$;
\item the Gaussian row on $-2^k$ removes all remaining node degrees below $n$ and returns the
      leading degree $n$.
\end{enumerate}
The exponent $n+k$ in the inner power is therefore forced: $k$ degrees are consumed by the
binary block, leaving exactly $n$ degrees for the geometric extractor.  The formula is a
hybrid of automatic-sequence cancellation, ordinary exponential generating functions, and
finite $q$-calculus; it is not a terminating basic-hypergeometric summation in the usual
${}_r\phi_s$ sense \cite{gasper2004,dlmf17}.

\section{Binary Taylor reduction and the global series}
\label{sec:binary}

\subsection{The finite reduction polynomial}

The inverse-dyadic moments assemble into the polynomial
\begin{equation}
 \mathcal T_m(y)=
 \sum_{j=0}^{m}
 2^{\binom{j+1}{2}-\binom{m-j}{2}}
 \frac{d_{m-j}}{(m-j)!}\frac{y^j}{j!}.                    \label{eq:T-def-local}
\end{equation}
The signed global extension satisfies the exact Taylor-block identity
\begin{equation}
 \mathcal F(2^{-m}+y)=-\mathcal F(y)+\mathcal T_m(y),
 \qquad 0\le y<2^{-m}.                                    \label{eq:T-reduction-local}
\end{equation}
This is the elementary step behind both the terminating dyadic recursion and the infinite
binary telescope.

For $x\ge0$ set
\begin{equation}
 p_m(x)=\lfloor2^mx\rfloor,
 \qquad
 y_m(x)=x-\frac{p_m(x)}{2^m},
 \qquad 0\le y_m(x)<2^{-m},                               \label{eq:prefix-tail-local}
\end{equation}
and put
\begin{equation}
 p^*_{m-1}(x)=
 \begin{cases}
 \lfloor x/2\rfloor,&m=0,\\
 p_{m-1}(x),&m\ge1.
 \end{cases}                                               \label{eq:previous-prefix-local}
\end{equation}
Define
\begin{equation}
 \alpha_m(x)=(-1)^{\wt(p_m(x))}
 \bigl(2p^*_{m-1}(x)-p_m(x)\bigr).                         \label{eq:alpha-local}
\end{equation}
This coefficient is $0$ unless the newly exposed binary digit is $1$, in which case it is a
sign.

\begin{theorem}[Binary reduction series]
\label{thm:binary-series-local}
For every $x\ge0$,
\begin{equation}
 \boxed{
 \mathcal F(x)=\sum_{m=0}^{\infty}
 \alpha_m(x)\mathcal T_m(y_m(x)).}                        \label{eq:binary-series-local}
\end{equation}
The series converges absolutely, and its partial sum through $m=N$ has the uniform remainder
bound
\begin{equation}
 \left|\mathcal F(x)-\sum_{m=0}^{N}
 \alpha_m(x)\mathcal T_m(y_m(x))\right|\le2^{1-N}
 \qquad(N\ge1).                                           \label{eq:binary-tail-local}
\end{equation}
\end{theorem}

The proof is an exact residual telescope.  If
\(
 R_m(x)=(-1)^{\wt(p_m(x))}\mathcal F(y_m(x)),
\)
then
\[
 \alpha_m(x)\mathcal T_m(y_m(x))=R_{m-1}(x)-R_m(x)
 \quad(m\ge1),
\]
with a separately correct scale-zero term.  Since $0\le y_m<2^{-m}$ and
$|\mathcal F(y)|\le2y$ near the origin, $|R_m|\le2^{1-m}$.

\subsection{The sparse digit form on the unit interval}

For $0\le x<1$, write the floor-selected binary expansion
\begin{equation}
 x=0.\delta_1\delta_2\delta_3\cdots
 =\sum_{m\ge1}\delta_m2^{-m},
 \qquad \delta_m\in\{0,1\},                               \label{eq:binary-digits}
\end{equation}
where $p_m$ is the integer represented by the first $m$ digits.  Let
\(
 s_m=\delta_1+\cdots+\delta_m
\)
and let
\begin{equation}
 t_m=\sum_{\ell>m}\delta_\ell2^{-\ell}=y_m(x)             \label{eq:binary-tail-digits}
\end{equation}
be the unscaled tail.  If $\delta_m=1$, then
$p_m=2p_{m-1}+1$ and
\[
 \alpha_m=(-1)^{s_m}(-1)=(-1)^{s_{m-1}};
\]
if $\delta_m=0$, then $\alpha_m=0$.

\begin{corollary}[Digit-sparse series]
\label{cor:digit-sparse}
For $0\le x<1$,
\begin{equation}
 \boxed{
 F(x)=\sum_{\substack{m\ge1\\\delta_m=1}}
 (-1)^{s_{m-1}}\mathcal T_m(t_m).}                         \label{eq:digit-sparse}
\end{equation}
At a dyadic rational the floor convention chooses the terminating binary expansion, so the
sum is finite and contains exactly one term for each set bit of the numerator.
\end{corollary}

This is the most direct connection between the fast bit-recursive evaluator and the arbitrary
argument series: the finite recursion is not a separate algorithm pasted onto the infinite
formula; it is literally the same telescope when the binary tail eventually becomes zero.

\subsection[Replacing each Taylor coefficient by a finite q-identity]{Replacing each Taylor coefficient by a finite $q$-identity}

For $Q\in\C$ define the finite inverse-block numerator
\begin{equation}
 \mathcal N_n^{(Q)}=
 \sum_{k=0}^{n}
 \frac{\qbinom nk{1/2}}{4^{\binom k2}(n+k)!}
 \sum_{r=0}^{2^k-1}\tm_r(r-2^k+Q)^{n+k}.                  \label{eq:NnQ-local}
\end{equation}
The exact one-block $q$-formula is equivalent to
\begin{equation}
 \frac{\mathcal N_n^{(Q)}}
 {2^{\binom n2}(1/2;1/2)_n}
 =\frac{2^nd_n}{n!},                                      \label{eq:NnQ-normalization-local}
\end{equation}
which is independent of $Q$.  Insert these finite identities into the polynomial
\eqref{eq:T-def-local}:
\begin{equation}
 \mathcal Q_m^{(Q)}(y)=
 2^{-\binom{m+1}{2}}
 \sum_{n=0}^{m}
 \frac{(2^{m+1}y)^{m-n}}{(m-n)!}
 \frac{\mathcal N_n^{(Q)}}
 {2^{\binom n2}(1/2;1/2)_n}.                              \label{eq:QmQ-local}
\end{equation}

\begin{theorem}[$q$-Taylor identity and global $q$-series]
\label{thm:q-global-local}
For every $m$, $y$, and $Q\in\C$,
\begin{equation}
 \mathcal Q_m^{(Q)}(y)=\mathcal T_m(y).                   \label{eq:Qm-equals-Tm-local}
\end{equation}
Consequently, for fixed $Q$ and every $x\ge0$,
\begin{equation}
 \boxed{
 \mathcal F(x)=\sum_{m=0}^{\infty}
 \alpha_m(x)\mathcal Q_m^{(Q)}(y_m(x)),}                 \label{eq:global-q-series-local}
\end{equation}
with absolute convergence.
\end{theorem}

The key word is \emph{termwise}: every finite polynomial
$\mathcal Q_m^{(Q)}$ already equals $\mathcal T_m$ and is already independent of $Q$ before
the outer infinite sum is considered.  The $q$-symbols do not create the convergence of the
outer series.  They merely provide an alternative finite presentation of each Taylor block.

\subsection{The dyadic triangle of finite representations}

Suppose $x=m/2^s\in[0,1]$ is written with a terminating binary expansion.  Then three different
finite mechanisms evaluate the same rational number:
\begin{enumerate}
\item \eqref{eq:input-moment-dyadic} sums the Thue--Morse prefix $0\le h<m$ and applies the
      even-moment correction in one shot;
\item \eqref{eq:input-q-dyadic} uses one geometric row $0\le k\le s$ and two nested
      cancellations;
\item \eqref{eq:digit-sparse} walks only through the set bits of $m$, using inverse-dyadic
      Taylor data.
\end{enumerate}
The finite sums have no natural term-by-term identification.  Their equality is mediated by
the master coefficient \eqref{eq:master-coefficient} and by the exact Taylor reduction
\eqref{eq:T-reduction-local}.  A fourth finite representation, obtained from the discrete
limit, appears in the next section.

\section[Discrete q-limits, exact stabilization, and extrapolation]{Discrete $q$-limits, exact stabilization, and extrapolation}
\label{sec:discrete}

\subsection{Shifted splines and the finite row}

For $p\ge1$, $Q\in\C$, and $x\in\R$, define
\begin{equation}
 S_p^{(Q)}(x)=
 \frac{(-1)^p}{2^{\binom p2}p!}
 \sum_{r=0}^{N_p(x)-1}\tm_r(r-2^px+Q)^p.                  \label{eq:shifted-spline-local}
\end{equation}
At $Q=1/2$ this is the centered spline $S_p$.  The formal development proves the exact
fixed-cutoff Taylor relation and the uniform bound
\begin{equation}
 \bigl|S_p^{(Q)}(x)-S_p(x)\bigr|
 \le2^{-(p-1)}\bigl(\e^{|Q-1/2|}-1\bigr).                 \label{eq:shifted-spline-bound-local}
\end{equation}
Together with \eqref{eq:centered-global-error-input},
\begin{equation}
 \bigl|S_p^{(Q)}(x)-\mathcal F(x)\bigr|
 \le(2\e^{|Q-1/2|}-1)2^{-p}.                              \label{eq:shifted-total-error}
\end{equation}

The discrete $q$-approximant is
\begin{align}
 \mathcal D_n^{(Q)}(x)
 &=\frac1{2^{n^2}(1/2;1/2)_n}
 \sum_{k=0}^{n}
 \frac{\qbinom nk{1/2}}{4^{\binom k2}(n+k)!}             \notag\\
 &\qquad\times
 \sum_{r=0}^{N_{n+k}(x)-1}\tm_r
       (r-2^{n+k}x+Q)^{n+k}.                              \label{eq:DnQ-local}
\end{align}
For every fixed $Q$ and every real $x$, the Lean development proves
\begin{equation}
 \mathcal D_n^{(Q)}(x)\longrightarrow\mathcal F(x).       \label{eq:DnQ-limit-local}
\end{equation}
The notation strongly resembles the exact dyadic formula
\eqref{eq:input-q-dyadic}.  The following result makes the resemblance literal.

\subsection{At dyadic inputs the limit stabilizes exactly}

\begin{theorem}[Finite stabilization on the dyadic grid]
\label{thm:dyadic-stabilization}
Let $x=m/2^s$ with $m,s\in\N$.  Then, for every $n\ge s$ and every $Q\in\C$,
\begin{equation}
 \boxed{\mathcal D_n^{(Q)}(x)=\mathcal F(x).}              \label{eq:dyadic-stabilization}
\end{equation}
If $0\le x\le1$, the right side is $F(x)$.
\end{theorem}

\begin{proof}
For $n\ge s$ put $M=m2^{n-s}$, so that $x=M/2^n$.  For every $0\le k\le n$,
\[
 2^{n+k}x=M2^k\in\N,
 \qquad
 N_{n+k}(x)=\left\lfloor M2^k+\frac12\right\rfloor=M2^k.
\]
Therefore \eqref{eq:DnQ-local} becomes exactly the dyadic formula
\eqref{eq:input-q-dyadic} with numerator $M$ and denominator exponent $n$.  That formula is
valid for unreduced dyadic representations and for every complex translation $Q$.
\end{proof}

This is the sharpest connection between the exact and limiting formulae: the discrete formula
extends the exact dyadic row to arbitrary $x$, and on a dyadic grid it needs no limiting
process once the row depth reaches the denominator depth.

\begin{corollary}[Exact finite shifted-spline extrapolation]
\label{cor:exact-shifted-extrapolation}
For $x=m/2^s$, $n\ge s$, and every $Q\in\C$,
\begin{equation}
 \mathcal F(x)=\sum_{j=0}^{n}\omega_{n,j}
 S_{2n-j}^{(Q)}(x),                                       \label{eq:exact-shifted-extrapolation}
\end{equation}
where the weights $\omega_{n,j}$ are defined in
\eqref{eq:omega-local} below.
\end{corollary}

Thus every dyadic value admits an exact representation as a finite signed combination of
consecutive spline degrees, all evaluated at the same point.  Each shifted spline may depend
strongly on $Q$, but the complete combination does not.

\subsection{The row polynomial}

Reindexing \eqref{eq:DnQ-local} by $j=n-k$ gives the exact identity
\begin{equation}
 \mathcal D_n^{(Q)}(x)=
 \sum_{j=0}^{n}\omega_{n,j}S_{2n-j}^{(Q)}(x),              \label{eq:D-toeplitz-local}
\end{equation}
with
\begin{equation}
 \omega_{n,j}=
 \frac{(-1)^j\qbinom nj{1/2}2^{-\binom{j+1}{2}}}
 {(1/2;1/2)_n}.                                           \label{eq:omega-local}
\end{equation}
The full row is encoded by one normalized $q$-Pochhammer polynomial.

\begin{proposition}[Generating polynomial of the Toeplitz row]
\label{prop:Wn}
Define
\begin{equation}
 W_n(z)=\sum_{j=0}^{n}\omega_{n,j}z^j.                    \label{eq:Wn-def}
\end{equation}
Then
\begin{equation}
 \boxed{
 W_n(z)=\frac{(z/2;1/2)_n}{(1/2;1/2)_n}
 =\prod_{\ell=1}^{n}\frac{1-z/2^\ell}{1-1/2^\ell}.}      \label{eq:Wn-product}
\end{equation}
Consequently
\begin{align}
 W_n(1)&=1,                                                \label{eq:Wn-one}\\
 W_n(2^r)&=0\qquad(1\le r\le n),                         \label{eq:Wn-roots}\\
 \sum_{j=0}^{n}|\omega_{n,j}|
 &=W_n(-1)=\frac{(-1/2;1/2)_n}{(1/2;1/2)_n},              \label{eq:Wn-variation}\\
 \sum_{j=0}^{n}|\omega_{n,j}|2^j
 &=W_n(-2)=\frac{(-1;1/2)_n}{(1/2;1/2)_n}.                \label{eq:Wn-weighted-variation}
\end{align}
\end{proposition}

\begin{proof}
Apply the finite $q$-binomial theorem to $(z/2;1/2)_n$.  Its coefficient of $z^j$ is
\[
 (-1)^j2^{-\binom j2}2^{-j}\qbinom nj{1/2}
 =(-1)^j2^{-\binom{j+1}{2}}\qbinom nj{1/2},
\]
which gives \eqref{eq:Wn-product} after division by $(1/2;1/2)_n$.  The evaluations at
$1,2^r,-1,-2$ follow from the product and from the alternating signs of the coefficients.
\end{proof}

The familiar row-sum identity and bounded variation are therefore only two evaluations of a
single polynomial.  The roots $2,4,\ldots,2^n$ contain additional information that the basic
Toeplitz proof does not need.

\subsection[A q-Pochhammer filter in the degree variable]{A $q$-Pochhammer filter in the degree variable}

Fix $x$ and $Q$ and regard $p\mapsto S_p^{(Q)}(x)$ as a sequence.  Let the backward shift
$\shiftop$ act by
\(
 \shiftop S_p^{(Q)}=S_{p-1}^{(Q)}.
\)
Then \eqref{eq:D-toeplitz-local} and \eqref{eq:Wn-product} can be written as
\begin{equation}
 \boxed{
 \mathcal D_n^{(Q)}(x)
 =W_n(\shiftop)S_{2n}^{(Q)}(x)
 =\prod_{\ell=1}^{n}
   \frac{1-2^{-\ell}\shiftop}{1-2^{-\ell}}
   S_{2n}^{(Q)}(x).}                                      \label{eq:q-filter-operator}
\end{equation}
Each factor preserves a constant sequence because its value at $\shiftop=1$ is $1$.
Meanwhile the zero $W_n(2^r)=0$ means that it removes a geometric error mode with ratio
$2^{-r}$ in the degree.

\begin{proposition}[$q$-Richardson exactness]
\label{prop:q-richardson}
Suppose a model sequence has the form
\begin{equation}
 u_p=L+\sum_{r=1}^{n}a_r2^{-rp}.                           \label{eq:model-errors}
\end{equation}
Then
\begin{equation}
 \sum_{j=0}^{n}\omega_{n,j}u_{2n-j}=L.                    \label{eq:q-richardson}
\end{equation}
\end{proposition}

\begin{proof}
The constant contributes $L W_n(1)=L$.  The $r$th error contributes
\[
 a_r2^{-2nr}\sum_{j=0}^{n}\omega_{n,j}2^{rj}
 =a_r2^{-2nr}W_n(2^r)=0.
\]
\end{proof}

No claim is needed that the actual spline error has a finite expansion of the form
\eqref{eq:model-errors}.  The proposition explains the algebraic design of the row: it is a
geometric-grid analogue of Richardson extrapolation, with the normalized $q$-Pochhammer
polynomial as its error-annihilating filter.

\subsection[A quantitative four-to-the-minus-n convergence bound]{A quantitative $4^{-n}$ convergence bound}

The bounded-variation proof of \eqref{eq:DnQ-limit-local} uses only that all sampled degrees
$2n-j$ tend to infinity.  Retaining their actual sizes gives a stronger quantitative
consequence.

\begin{theorem}[Explicit discrete-row error]
\label{thm:D-error}
For $n\ge1$, $Q\in\C$, and every $x\in\R$, let
\begin{equation}
 H_n=\frac{(-1;1/2)_n}{(1/2;1/2)_n}.                      \label{eq:Hn-def}
\end{equation}
Then
\begin{equation}
 \boxed{
 |\mathcal D_n^{(Q)}(x)-\mathcal F(x)|
 \le H_n\bigl(2\e^{|Q-1/2|}-1\bigr)4^{-n}.}              \label{eq:D-error}
\end{equation}
Moreover $H_n\le64$, so the convergence is uniformly $\bigO_Q(4^{-n})$ on the whole real
line.  At the centered shift,
\begin{equation}
 |\mathcal D_n^{(1/2)}(x)-\mathcal F(x)|\le H_n4^{-n}.     \label{eq:D-centered-error}
\end{equation}
\end{theorem}

\begin{proof}
Apply \eqref{eq:shifted-total-error} with $p=2n-j$ in
\eqref{eq:D-toeplitz-local}:
\begin{align*}
 |\mathcal D_n^{(Q)}(x)-\mathcal F(x)|
 &\le\bigl(2\e^{|Q-1/2|}-1\bigr)
 \sum_{j=0}^{n}|\omega_{n,j}|2^{-(2n-j)}\\
 &=\bigl(2\e^{|Q-1/2|}-1\bigr)4^{-n}
 \sum_{j=0}^{n}|\omega_{n,j}|2^j.
\end{align*}
Use \eqref{eq:Wn-weighted-variation}.  Finally
\[
 H_n
 =\frac{2}{1-2^{-n}}
   \frac{(-1/2;1/2)_{n-1}}{(1/2;1/2)_{n-1}}
 \le4\cdot16=64,
\]
where the last factor is the uniform variation bound proved in the formal development.
\end{proof}

The same argument isolates the finite translation dependence.

\begin{corollary}[Decay of the shift dependence]
\label{cor:D-Q-decay}
For every $x\in\R$ and $n\ge1$,
\begin{equation}
 |\mathcal D_n^{(Q)}(x)-\mathcal D_n^{(1/2)}(x)|
 \le2H_n\bigl(\e^{|Q-1/2|}-1\bigr)4^{-n}.                \label{eq:D-Q-decay}
\end{equation}
\end{corollary}

Thus the finite rows can depend dramatically on $Q$ at small $n$, but the dependence is
uniformly suppressed at a geometric rate.  This is analytically different from
\cref{prop:finite-Q-invariance}, where the finite dyadic expression is exactly constant in
$Q$.

\subsection[A synchronized nearest-dyadic q-limit]{A synchronized nearest-dyadic $q$-limit}
\label{subsec:synchronized}

There is a simpler way to turn the exact dyadic formula into a limit at arbitrary argument.
For $x\ge0$ define the nearest nonnegative dyadic numerator
\begin{equation}
 M_n(x)=\max\!\left(0,\left\lfloor2^nx+\frac12\right\rfloor\right)=N_n(x),
 \qquad
 \widetilde x_n=\frac{M_n(x)}{2^n}.                        \label{eq:nearest-dyadic}
\end{equation}
Use this one numerator at every inner scale:
\begin{align}
 \widetilde{\mathcal D}_n^{(Q)}(x)
 &=\frac1{2^{n^2}(1/2;1/2)_n}
 \sum_{k=0}^{n}
 \frac{\qbinom nk{1/2}}{4^{\binom k2}(n+k)!}             \notag\\
 &\qquad\times
 \sum_{r=0}^{M_n(x)2^k-1}\tm_r
       (r-M_n(x)2^k+Q)^{n+k}.                             \label{eq:synchronized-D}
\end{align}
Unlike \eqref{eq:DnQ-local}, every finite row is now an exact dyadic formula.

\begin{theorem}[Synchronized exact-grid representation]
\label{thm:synchronized}
For every $n$, $Q\in\C$, and $x\ge0$,
\begin{equation}
 \boxed{
 \widetilde{\mathcal D}_n^{(Q)}(x)
 =\mathcal F\!\left(\frac{M_n(x)}{2^n}\right).}           \label{eq:synchronized-exact}
\end{equation}
In particular, the left side is exactly independent of $Q$.  On $0\le x\le1$,
\begin{equation}
 \left|\widetilde{\mathcal D}_n^{(Q)}(x)-F(x)\right|
 \le2^{-n},                                                \label{eq:synchronized-error}
\end{equation}
and hence
\begin{equation}
 F(x)=\lim_{n\to\infty}\widetilde{\mathcal D}_n^{(Q)}(x). \label{eq:synchronized-limit}
\end{equation}
\end{theorem}

\begin{proof}
Equation \eqref{eq:synchronized-exact} is
\eqref{eq:input-q-dyadic} with numerator $M_n(x)$.  For $x\in[0,1]$,
$|\widetilde x_n-x|\le2^{-n-1}$ and the bounded Fabius function is $2$-Lipschitz, giving
\eqref{eq:synchronized-error}.
\end{proof}

The same construction applied to \eqref{eq:input-moment-dyadic} gives a purely
moment/Thue--Morse limit:
\begin{align}
 F(x)=\lim_{n\to\infty}
 &\frac{2^{-\binom{n+1}{2}}}{n!}
 \sum_{h=0}^{M_n(x)-1}\tm_h                              \notag\\
 &\times\sum_{k=0}^{\lfloor n/2\rfloor}
 \binom n{2k}c_k(2M_n(x)-2h-1)^{n-2k},                   \label{eq:synchronized-moment-limit}
\end{align}
again with error at most $2^{-n}$ on $[0,1]$.  This representation contains no $q$-symbols at
all: it is exact dyadic arithmetic followed by nearest-grid convergence.

The synchronized limit is conceptually simple but converges at the basic grid rate
$\bigO(2^{-n})$.  The asynchronous row \eqref{eq:DnQ-local} is adapted to the natural spline
degree $n+k$ at each scale and, by \cref{thm:D-error}, achieves the stronger uniform bound
$\bigO_Q(4^{-n})$.  The normalized $q$-Pochhammer row acts as an acceleration filter.

\subsection{What fails away from a dyadic point}

Let
\begin{equation}
 \theta_n(x)=2^nx-M_n(x),
 \qquad -\frac12\le\theta_n(x)<\frac12.                  \label{eq:theta-n}
\end{equation}
Then for every $k\ge0$,
\begin{equation}
 N_{n+k}(x)
 =M_n(x)2^k+\left\lfloor2^k\theta_n(x)+\frac12\right\rfloor, \label{eq:cutoff-defect}
\end{equation}
and the shifted coordinate in the inner power is
\begin{equation}
 r-2^{n+k}x+Q
 =r-M_n(x)2^k+\bigl(Q-2^k\theta_n(x)\bigr).               \label{eq:shift-defect}
\end{equation}
Thus the source discrete row differs from the synchronized exact dyadic row in exactly two
ways:
\begin{enumerate}
\item the complete block of length $M_n2^k$ acquires a boundary defect
      $\lfloor2^k\theta_n+1/2\rfloor$;
\item the common translation $Q$ is replaced by the scale-dependent translation
      $Q-2^k\theta_n$.
\end{enumerate}
At a resolved dyadic point, $\theta_n=0$, so both defects disappear simultaneously; this is
another proof of \cref{thm:dyadic-stabilization}.  Away from the dyadic grid the shifts vary
with $k$, so the finite common-translation invariance of
\cref{prop:finite-Q-invariance} cannot be applied across the row.

\subsection{Two infinite constructions that share a limit but not terms}

The global binary $q$-series \eqref{eq:global-q-series-local} and the discrete limit
\eqref{eq:DnQ-limit-local} both evaluate $\mathcal F$, but their relation is not that of an
infinite series and its partial sums.

\begin{table}[ht]
\centering
\small
\begin{tabularx}{\textwidth}{@{}p{0.22\textwidth}XX@{}}
\toprule
Feature & Binary $q$-series & Discrete $q$-limit \\
\midrule
Outer index & binary scale $m$ & approximation row $n$ \\
Finite inner object & $\mathcal Q_m^{(Q)}=\mathcal T_m$ & shifted splines $S_{2n-j}^{(Q)}$ \\
Shift dependence & absent in every summand & generally present in every finite row \\
Convergence source & exact residual telescope & Toeplitz summability of spline convergence \\
At dyadic $x$ & outer series terminates at the last set bit & rows stabilize exactly once $n$ reaches the denominator depth \\
Natural error bound & $2^{1-N}$ after scale $N$ & $H_n(2\e^{|Q-1/2|}-1)4^{-n}$ \\
\bottomrule
\end{tabularx}
\caption{The two arbitrary-argument $q$ constructions.}
\label{tab:two-infinite}
\end{table}

\begin{warning}
It is therefore incorrect to call $\mathcal D_n^{(Q)}$ the $n$th partial sum of the global
$q$-binomial series.  The two families use the same finite half-base identities, but they
insert them into different analytic frameworks.
\end{warning}

\subsection{The first rows}

For reference, the first normalized Toeplitz rows are
\begin{align*}
 (\omega_{1,j})_{j=0}^{1}&=(2,-1),\\
 (\omega_{2,j})_{j=0}^{2}&=(8/3,-2,1/3),\\
 (\omega_{3,j})_{j=0}^{3}&=(64/21,-8/3,2/3,-1/21),\\
 (\omega_{4,j})_{j=0}^{4}&=(1024/315,-64/21,8/9,-2/21,1/315).
\end{align*}
Their ordinary and $2^j$-weighted variations are shown in
\cref{tab:row-constants}.

\begin{table}[ht]
\centering
\small
\begin{tabular}{@{}c c c@{}}
\toprule
$n$ & $\sum_j|\omega_{n,j}|$ & $H_n=\sum_j|\omega_{n,j}|2^j$ \\
\midrule
1 & $3$ & $4$ \\
2 & $5$ & $8$ \\
3 & $45/7$ & $80/7$ \\
4 & $51/7$ & $96/7$ \\
5 & $1683/217$ & $3264/217$ \\
\bottomrule
\end{tabular}
\caption{Two evaluations of the row polynomial: $W_n(-1)$ and $W_n(-2)$.}
\label{tab:row-constants}
\end{table}

\section{A unified dyadic identity and a worked example}
\label{sec:worked}

\subsection{The complete equality chain}

Let $x=m/2^s\in[0,1]$, use the terminating binary digits
$x=0.\delta_1\cdots\delta_s$, and let $N\ge s$.  Combining the preceding results gives the
following chain of exact finite expressions:
\begin{align}
 F(x)
 &=2^{-\binom s2}[X^s]\bigl(B_m(X)\mathscr A(X)\bigr)     \label{eq:chain-coefficient}\\
 &=\frac{2^{-\binom{s+1}{2}}}{s!}
   \sum_{h=0}^{m-1}\tm_h
   \sum_{k=0}^{\lfloor s/2\rfloor}
   \binom s{2k}c_k(2m-2h-1)^{s-2k}                       \label{eq:chain-moment}\\
 &=\frac{1}{2^{s^2}(1/2;1/2)_s}
   \sum_{k=0}^{s}\frac{\qbinom sk{1/2}}{4^{\binom k2}(s+k)!}
   \sum_{r=0}^{m2^k-1}\tm_r(r-m2^k+Q)^{s+k}              \label{eq:chain-q}\\
 &=\sum_{k=0}^{\lfloor s/2\rfloor}
   \frac{c_k}{(2k)!2^{k(2s-2k+1)}}
   S_{s-2k}\!\left(\frac{m}{2^{s-2k}}\right)             \label{eq:chain-moment-spline}\\
 &=\sum_{\substack{1\le j\le s\\\delta_j=1}}
   (-1)^{\delta_1+\cdots+\delta_{j-1}}
   \mathcal T_j\!\left(\sum_{\ell>j}\delta_\ell2^{-\ell}\right) \label{eq:chain-binary}\\
 &=\sum_{j=0}^{N}\omega_{N,j}S_{2N-j}^{(Q)}(x).           \label{eq:chain-q-spline}
\end{align}
The letter $k$ has unrelated local meanings in different lines; this is unavoidable because
one line indexes even moments and another indexes geometric scales.  The equalities are not
termwise transformations.  They are different factorizations of the same exact value:
\begin{itemize}
\item \eqref{eq:chain-coefficient} is the common scalar invariant;
\item \eqref{eq:chain-moment} expands the probability law in centered moments;
\item \eqref{eq:chain-q} extracts the same coefficient by geometric interpolation;
\item \eqref{eq:chain-moment-spline} corrects one matched spline through even moments;
\item \eqref{eq:chain-binary} telescopes through the set bits;
\item \eqref{eq:chain-q-spline} extrapolates consecutive spline degrees at the same point.
\end{itemize}

\subsection[The value at 5/16]{The value at $5/16$}

Take
\[
 x=\frac5{16}=0.0101_2,
 \qquad
 F(x)=\frac{305857}{2073600}
 \approx0.1475004822530864.                                \label{eq:F-five-sixteen}
\]
This one value illustrates every bridge above.

\subsubsection{Master coefficient and moment formula}

The prefix polynomial is
\begin{equation}
 B_5(X)=\e^{4X}-\e^{3X}-\e^{2X}+\e^X-1.                  \label{eq:B5}
\end{equation}
Using
\(
 d_0=1,d_1=1/2,d_2=5/18,d_3=1/6,d_4=143/1350
\)
in $\mathscr A$ gives
\begin{equation}
 [X^4](B_5\mathscr A)=\frac{305857}{32400}.               \label{eq:B5-coefficient}
\end{equation}
Since $2^{-\binom42}=1/64$, \eqref{eq:master-coefficient} gives
\eqref{eq:F-five-sixteen}.  Equivalently, with
$c_1=1/9$ and $c_2=19/675$,
\begin{equation}
 F(5/16)=\frac{2^{-10}}{4!}
 \sum_{h=0}^{4}\tm_h
 \left[t_h^4+6c_1t_h^2+c_2\right],
 \qquad t_h=9-2h.                                         \label{eq:five-sixteen-moment}
\end{equation}

\subsubsection{Moment correction of the matched spline}

The three spline values required by \eqref{eq:moment-corrected-spline} are
\begin{equation}
 S_4(5/16)=\frac{1205}{8192},
 \qquad
 S_2(5/4)=\frac{15}{16},
 \qquad
 S_0(5)=-1.                                                \label{eq:five-sixteen-splines}
\end{equation}
Hence
\begin{align}
 F(5/16)
 &=S_4(5/16)+\frac1{2304}S_2(5/4)
   +\frac{19}{16588800}S_0(5)                             \notag\\
 &=\frac{1205}{8192}+\frac5{12288}-\frac{19}{16588800}
 =\frac{305857}{2073600}.                                 \label{eq:five-sixteen-correction}
\end{align}
The matched degree-four spline is already close, but the exact discrepancy is resolved by two
universal lower-degree corrections.

\subsubsection{The two set bits}

The binary digits equal $1$ only at positions $2$ and $4$.  The reduction polynomial at the
first position is
\begin{equation}
 \mathcal T_2(y)=\frac5{72}+y+4y^2,                        \label{eq:T2-explicit}
\end{equation}
and
\(
 \mathcal T_4(0)=F(1/16)=143/2073600.
\)
The signs are $+$ at position $2$ and $-$ at position $4$, so
\begin{align}
 F(5/16)
 &=\mathcal T_2(1/16)-\mathcal T_4(0)                     \notag\\
 &=\frac{85}{576}-\frac{143}{2073600}
 =\frac{305857}{2073600}.                                 \label{eq:five-sixteen-binary}
\end{align}
The prefix formula sums five Thue--Morse terms; the binary formula needs only the two set bits.

\subsubsection[Exact q-extrapolation of centered splines]{Exact $q$-extrapolation of centered splines}

At row $N=4$, \eqref{eq:chain-q-spline} reads
\begin{align}
 F(5/16)
 ={}&\frac{1024}{315}S_8(5/16)
 -\frac{64}{21}S_7(5/16)
 +\frac89S_6(5/16)                                      \notag\\
 &-\frac2{21}S_5(5/16)
 +\frac1{315}S_4(5/16).                                  \label{eq:five-sixteen-q-extrap}
\end{align}
The individual spline values are
\begin{equation}
 \begin{split}
 S_4&=\frac{1205}{8192},\qquad
 S_5=\frac{289799}{1966080},\qquad
 S_6=\frac{13917509}{94371840},\\
 S_7&=\frac{74236247}{503316480},\qquad
 S_8=\frac{395939357}{2684354560},
 \end{split}                                               \label{eq:five-sixteen-S4-S8}
\end{equation}
all evaluated at $5/16$.  None equals $F(5/16)$, but the $q$-Pochhammer filter combines them
exactly.

\subsubsection{Finite shift dependence before stabilization}

The first discrete rows illustrate the difference between exact finite translation invariance
and asymptotic translation independence.

\begin{table}[ht]
\centering
\small
\begin{tabular}{@{}c c c c@{}}
\toprule
$n$ & $\mathcal D_n^{(0)}(5/16)$ & $\mathcal D_n^{(1/2)}(5/16)$ & $\mathcal D_n^{(3)}(5/16)$ \\
\midrule
1 & $0.156250000$ & $0.156250000$ & $3.906250000$ \\
2 & $0.151041667$ & $0.147460938$ & $-0.531250000$ \\
3 & $0.147505374$ & $0.147501983$ & $0.125788484$ \\
4 & $0.147500482$ & $0.147500482$ & $0.147500482$ \\
\bottomrule
\end{tabular}
\caption{The row becomes exactly equal to $F(5/16)$ at $n=4$, independently of $Q$.}
\label{tab:five-sixteen-D}
\end{table}

The exact dyadic $q$-formula itself is independent of $Q$ already at row four, although its
outer summands vary substantially.  Their decimal contributions are:

\begin{table}[ht]
\centering
\small
\begin{tabular}{@{}c r r r@{}}
\toprule
scale $k$ & $Q=0$ & $Q=1/2$ & $Q=3$ \\
\midrule
0 & $\phantom{-}0.000626240$ & $\phantom{-}0.000466967$ & $\phantom{-}0.000000000$ \\
1 & $-0.016345021$ & $-0.014038037$ & $-0.005184307$ \\
2 & $\phantom{-}0.141657926$ & $\phantom{-}0.131089095$ & $\phantom{-}0.084488329$ \\
3 & $-0.467452954$ & $-0.449506045$ & $-0.365062120$ \\
4 & $\phantom{-}0.489014292$ & $\phantom{-}0.479488502$ & $\phantom{-}0.433258580$ \\
\midrule
sum & $0.147500482$ & $0.147500482$ & $0.147500482$ \\
\bottomrule
\end{tabular}
\caption{Different translated decompositions of the same exact dyadic value.}
\label{tab:five-sixteen-q-terms}
\end{table}

\section{Computational and conceptual consequences}
\label{sec:consequences}

\subsection{Which representation is best for which task?}

The formulae are mathematically equivalent at dyadic inputs, but not computationally
interchangeable.

\begin{table}[ht]
\centering
\small
\begin{tabularx}{\textwidth}{@{}p{0.25\textwidth}p{0.31\textwidth}X@{}}
\toprule
Task & Recommended representation & Reason \\
\midrule
Prove rationality or denominator facts & master coefficient or moment formula & all arithmetic is finite and the universal moments are explicit rationals \\
Evaluate one dyadic with a sparse numerator & binary Taylor reduction & the number of reductions is the number of set bits \\
Evaluate many numerators at one depth & precompute $d_j$ or $c_k$, then use prefix recurrences & universal moment data are reused \\
Explain the $q$-coefficients & geometric divided difference and residue form & Gaussian coefficients are recognized as interpolation weights, not mysterious special-function decorations \\
Approximate arbitrary $x$ with a simple certified error & centered spline $S_p$ & direct global error $2^{-p}$ and one finite Thue--Morse sum \\
Accelerate arbitrary-argument spline convergence & discrete row $\mathcal D_n^{(Q)}$ & exact $q$-Pochhammer filter and uniform $\bigO_Q(4^{-n})$ bound \\
Use only exact dyadic arithmetic in a limiting scheme & synchronized row $\widetilde{\mathcal D}_n^{(Q)}$ or synchronized moment limit & every finite stage is an exact dyadic value and is translation-independent \\
Trace dependence on binary digits & digit-sparse series & one nonzero summand per exposed $1$-bit \\
\bottomrule
\end{tabularx}
\caption{Representation selection.  The $q$-formula is structurally illuminating but is
usually not the cheapest naive exact evaluator.}
\label{tab:representation-selection}
\end{table}

The exact $q$-formula contains inner blocks of length $m2^k$, so a direct implementation can
be expensive.  Its value is primarily structural: it exposes geometric interpolation,
translation invariance, and the extrapolation filter.  For exact numerical evaluation, the
bit recursion or the master coefficient with precomputed moments is generally preferable.

\subsection{A compact exact-arithmetic scheme}

Equation \eqref{eq:d-recurrence-companion} computes all half moments with rational arithmetic,
after which \eqref{eq:d-moment-prefix} evaluates a dyadic value.

\begin{algorithm}[Moment-prefix evaluator]
\label{alg:moment-prefix}
Given $m,n\in\N$:
\begin{enumerate}
\item Set $d_0=1$ and, for $1\le j\le n$, compute
\[
 d_j=\frac1{2^j-1}\sum_{r=1}^{j}\binom jr\frac{d_{j-r}}{r+1}.
\]
\item Form
\[
 V=\sum_{h=0}^{m-1}\tm_h
   \sum_{j=0}^{n}\binom njd_j(m-1-h)^{n-j}.
\]
\item Return $2^{-\binom n2}V/n!$.
\end{enumerate}
The result is $\mathcal F(m/2^n)$ and is $F(m/2^n)$ when $0\le m\le2^n$.
\end{algorithm}

A literal pseudocode version is included for clarity.

\begin{lstlisting}[style=pseudo,caption={Exact dyadic evaluation from half moments.}]
d[0] := 1
for j := 1,...,n do
    s := 0
    for r := 1,...,j do
        s := s + binomial(j,r) * d[j-r] / (r+1)
    end
    d[j] := s / (2^j - 1)
end

V := 0
for h := 0,...,m-1 do
    inner := 0
    for j := 0,...,n do
        inner := inner + binomial(n,j) * d[j] * (m-1-h)^(n-j)
    end
    V := V + (-1)^(binary_weight(h)) * inner
end
return V / (2^(n*(n-1)/2) * n!)
\end{lstlisting}

The centered-moment formula can be evaluated similarly.  Its advantage is parity: only
$\lfloor n/2\rfloor+1$ moments occur.  The half-moment formula aligns more directly with the
binary Taylor polynomials.

\subsection{Exact arithmetic checks implied by the connections}

The network of representations supplies a strong verification suite.  For chosen $m,n$ one
can independently check:
\begin{enumerate}
\item the master coefficient against the centered-moment formula;
\item the half-base $q$ expression at two or more translations $Q$;
\item the terminating binary reduction when $m/2^n\in[0,1]$;
\item the moment-corrected matched-spline identity;
\item the exact $q$-extrapolation row at any order $N\ge n$;
\item representation independence under $(m,n)\mapsto(2m,n+1)$.
\end{enumerate}
These checks stress different cancellation patterns.  Agreement is therefore more informative
than recomputing the same formula in a rearranged order.

\subsection{Endpoint and convention warnings}

Several small conventions are mathematically active.

\begin{warning}[The scale-zero term]
The global binary series begins at $m=0$.  It happens to vanish for $0\le x<1$, but it equals
$1$ at $x=1$ and controls the signed behavior on integer blocks.  A series beginning at
$m=1$ is valid only on the half-open unit interval.
\end{warning}

\begin{warning}[Integer exponents]
The exponent
$\binom{j+1}{2}-\binom{m-j}{2}$ in \eqref{eq:T-def-local} is an integer and may be negative.
Replacing integer subtraction by truncated natural subtraction changes the polynomial.
\end{warning}

\begin{warning}[Centered powers are signed]
Expressions such as $(r-m2^k+Q)^{n+k}$ live in a characteristic-zero scalar field.  The
quantity $r-m2^k$ is deliberately negative on much of the block.  Truncated subtraction in
$\N$ destroys the formula.
\end{warning}

\begin{warning}[Binary expansions at dyadics]
The floor prefixes $p_m=\lfloor2^mx\rfloor$ select the terminating expansion at a dyadic
rational.  Replacing it by the alternative expansion ending in infinitely many $1$-digits
changes the termwise digit series even though the represented real number is unchanged.
\end{warning}

\section{Relation to the Lean development}
\label{sec:formal-map}

The principal source theorems are formalized in the following modules.  Filenames are relative
to \repo.

\begingroup
\small
\setlength\LTleft{0pt}
\setlength\LTright{0pt}
\begin{longtable}{@{}p{0.34\textwidth}p{0.62\textwidth}@{}}
\toprule
Module & Role in the present article \\
\midrule
\endfirsthead
\toprule
Module & Role in the present article \\
\midrule
\endhead
\path{DyadicClosedForm.lean} & Exact Thue--Morse/even-moment formula for arbitrary dyadic numerators. \\
\path{HalfQBinomial.lean} & Finite $q$-Pochhammer and Gaussian-binomial calculus at $q=1/2$, including the interpolation annihilation row. \\
\path{ThueMorseExponential.lean} & Complete-block exponential factorization, Prouhet cancellation, and translation of inner powers. \\
\path{FabiusQBinomialFormula.lean} & Centered finite $q$-binomial formula, arbitrary numerators, and common-translation invariance. \\
\path{FabiusQBinomialTaylor.lean} & Identification of the finite $q$ expression with the Taylor reduction polynomial. \\
\path{FabiusQBinomialScalarFormula.lean} & Scalar normalizations and real/complex translated forms. \\
\path{FabiusRawQBinomialFormula.lean} and \path{FabiusRawQBinomialScalar.lean} & Raw inverse-dyadic forms and their scalar versions. \\
\path{FabiusBinaryReductionSeries.lean} & Exact binary residual telescope, absolute convergence, and uniform tail bound. \\
\path{FabiusParityPowerSeries.lean} & Corrected scale-zero parity-power presentation. \\
\path{FabiusGlobalQBinomialSeries.lean} & Termwise substitution of finite $q$-identities into the binary telescope. \\
\path{FabiusComplexShiftSpline.lean} & Shifted finite splines, fixed-cutoff Taylor expansion, and uniform shift bound. \\
\path{FabiusDiscreteLimitToeplitz.lean} & Row mass, exact total variation, uniform variation bound, and finite-row Toeplitz theorem. \\
\path{FabiusDiscreteLimitIntegration.lean} & Reindexing of the discrete formula and convergence to the signed global function. \\
\bottomrule
\end{longtable}
\endgroup

The following statements are direct paper-level corollaries or reorganizations of those
formalized results:
\begin{itemize}
\item the expectation and two contour representations
      \eqref{eq:expectation-Y}, \eqref{eq:cauchy-X}, and \eqref{eq:q-node-contour};
\item the composition and Bell-polynomial formulae for $d_n$;
\item the moment-corrected spline transform and its reciprocal
      \eqref{eq:moment-corrected-spline}--\eqref{eq:reciprocal-spline};
\item exact finite stabilization of the discrete row at dyadic arguments;
\item the row polynomial $W_n$, the shift-operator form, and the $q$-Richardson interpretation;
\item the explicit $\bigO_Q(4^{-n})$ bound \eqref{eq:D-error};
\item the synchronized nearest-dyadic $q$- and moment-limits.
\end{itemize}
They require no new unproved property of the Fabius function; each follows by finite algebra,
coefficient extraction, or the already formalized uniform spline estimates.  They are marked
here to distinguish conceptual reorganization from theorem names that already occur verbatim
in the repository.

\section{Conclusion}
\label{sec:conclusion}

The exact dyadic, global-series, and discrete-limit formulae are best understood not as a
collection of special tricks but as a small number of reusable operators acting in different
variables.

The product $B_m\mathscr A$ separates numerator arithmetic from the universal probability
law.  Expanding its coefficient in centered moments gives the classical Thue--Morse formula;
reading it as an expectation gives a probabilistic representation; taking a Cauchy integral
gives an analytic one.  Introducing the geometric node $z$ turns the same coefficient into
the leading term of a degree-$n$ polynomial.  The Gaussian row at base $1/2$ is then simply
the barycentric divided difference on $-1,-2,\ldots,-2^n$, and the $q$-Pochhammer symbol is the
node polynomial of that grid.

The centered finite spline sits directly inside the exact arithmetic.  Its matched-grid value
is the zeroth centered-moment term, and the remaining moments form an exact, invertible defect
correction.  A second correction uses the normalized $q$-Pochhammer polynomial in the spline
degree: it preserves constants, annihilates geometric error modes, and yields exact recovery
at dyadic points.  The resulting discrete row is therefore naturally viewed as a
$q$-Richardson extrapolator, not merely as an opaque Toeplitz average.

Finally, the two arbitrary-argument $q$ constructions have distinct logical roles.  The global
$q$-series is the binary residual telescope with an alternative finite formula inserted into
each summand.  The discrete limit is a finite extrapolation row whose shift dependence
disappears only asymptotically.  Their meeting point is the dyadic grid: there the binary
series terminates, the discrete row stabilizes, the synchronized row is exact, and all roads
return to the same master coefficient.

\appendix

\section{A compact identity sheet}
\label{app:identity-sheet}

For convenience, the principal transformations are collected here.

\subsection*{Probability and moments}
\begin{align*}
 Y&=\sum_{j\ge1}U_j/2^j,\qquad Z=2Y-1,\\
 \mathscr A(X)&=\E(\e^{XY})=\prod_{j\ge1}\frac{\e^{X/2^j}-1}{X/2^j},\\
 \mathscr A(X)&=\e^{X/2}\mathscr C(X/2),\qquad
 \mathscr C(T)=\sum_{k\ge0}c_kT^{2k}/(2k)!,\\
 d_n&=2^{-n}\sum_{k\le n/2}\binom n{2k}c_k,\\
 c_k&=\sum_{j=0}^{2k}(-1)^j2^j\binom{2k}{j}d_j.
\end{align*}

\subsection*{Master coefficient}
\begin{align*}
 B_m(X)&=\sum_{h<m}\tm_h\e^{(m-1-h)X},\\
 \mathcal F(m/2^n)&=2^{-\binom n2}[X^n](B_m\mathscr A),\\
 \mathcal F(m/2^n)&=\frac{2^{-\binom n2}}{n!}
   \sum_{h<m}\tm_h\E(m-1-h+Y)^n.
\end{align*}

\subsection*{Geometric extraction}
\begin{align*}
 w_{n,k}&=(-1)^k2^{-\binom k2}\qbinom nk{1/2},
 \qquad x_k=-2^k,\\
 \frac1{\mathfrak m_n}\sum_{k=0}^{n}w_{n,k}P(x_k)&=[z^n]P(z)
 \quad(\deg P\le n),\\
 \mathfrak m_n&=2^{\binom{n+1}{2}}(1/2;1/2)_n.
\end{align*}

\subsection*{Spline corrections}
\begin{align*}
 \mathcal F(m/2^n)
 &=\sum_{k\le n/2}
 \frac{c_k}{(2k)!2^{k(2n-2k+1)}}
 S_{n-2k}(m/2^{n-2k}),\\
 S_n(m/2^n)
 &=\sum_{k\le n/2}
 \frac{\gamma_k}{(2k)!2^{k(2n-2k+1)}}
 \mathcal F(m/2^{n-2k}),\\
 \sum_{k\ge0}\gamma_kT^{2k}/(2k)!&=1/\mathscr C(T).
\end{align*}

\subsection*{Discrete filter}
\begin{align*}
 \mathcal D_n^{(Q)}(x)&=\sum_{j=0}^{n}\omega_{n,j}S_{2n-j}^{(Q)}(x),\\
 W_n(z)=\sum_j\omega_{n,j}z^j
 &=\frac{(z/2;1/2)_n}{(1/2;1/2)_n},\\
 |\mathcal D_n^{(Q)}(x)-\mathcal F(x)|
 &\le\frac{(-1;1/2)_n}{(1/2;1/2)_n}
 (2\e^{|Q-1/2|}-1)4^{-n},\\
 x=m/2^s,\ n\ge s&\Longrightarrow
 \mathcal D_n^{(Q)}(x)=\mathcal F(x).
\end{align*}

\begin{thebibliography}{99}

\bibitem{fabius1966}
J.~Fabius,
\newblock A probabilistic example of a nowhere analytic $C^\infty$-function,
\newblock \emph{Zeitschrift fuer Wahrscheinlichkeitstheorie und Verwandte Gebiete}
  \textbf{5} (1966), 173--174.
\newblock DOI: \nolinkurl{10.1007/BF00536652}.

\bibitem{arias1982}
J.~Arias de Reyna,
\newblock Definicion y estudio de una funcion indefinidamente diferenciable de soporte
  compacto,
\newblock \emph{Revista de la Real Academia de Ciencias Exactas, Fisicas y Naturales de
  Madrid} \textbf{76} (1982), no.~1, 21--38.
\newblock English translation: \emph{An infinitely differentiable function with compact
  support: Definition and properties}, arXiv:\nolinkurl{1702.05442} (2017).

\bibitem{arias2018}
J.~Arias de Reyna,
\newblock Arithmetic of the Fabius function,
\newblock \emph{INTEGERS} \textbf{18} (2018), Article A51;
  arXiv:\nolinkurl{1702.06487}, version~3.

\bibitem{haugland2016}
J.~K.~Haugland,
\newblock Evaluating the Fabius function,
\newblock arXiv:\nolinkurl{1609.07999} (2016).

\bibitem{allouche2003}
J.-P.~Allouche and J.~Shallit,
\newblock \emph{Automatic Sequences: Theory, Applications, Generalizations},
\newblock Cambridge University Press, 2003.

\bibitem{andrews1976}
G.~E.~Andrews,
\newblock \emph{The Theory of Partitions},
\newblock Encyclopedia of Mathematics and its Applications~2,
  Addison--Wesley, 1976.

\bibitem{gasper2004}
G.~Gasper and M.~Rahman,
\newblock \emph{Basic Hypergeometric Series}, second edition,
\newblock Encyclopedia of Mathematics and its Applications~96,
  Cambridge University Press, 2004.

\bibitem{kac2002}
V.~Kac and P.~Cheung,
\newblock \emph{Quantum Calculus},
\newblock Universitext, Springer, 2002.

\bibitem{dlmf17}
NIST Digital Library of Mathematical Functions, Chapter~17,
\newblock \emph{$q$-Hypergeometric and Related Functions},
\newblock \url{https://dlmf.nist.gov/17}.

\bibitem{reshetnikovMSE}
V.~Reshetnikov,
\newblock A conjectured closed form for the Fabius function at dyadic rationals,
\newblock Mathematics Stack Exchange, question 3283519, 5 July 2019.

\bibitem{proveit}
V.~Reshetnikov,
\newblock \emph{ProveIt: Analysis/FabiusFunction},
\newblock Lean~4 formal development, repository snapshot
  \path{75dac857e96d858140e7e5520172d0c659bde98e}, 24 August 2026.
\newblock \url{https://github.com/VladimirReshetnikov/ProveIt/tree/75dac857e96d858140e7e5520172d0c659bde98e/Analysis/FabiusFunction}.

\bibitem{mainarticle}
V.~Reshetnikov,
\newblock \emph{The Fabius Function and Rvachev's Up-Function: Exact arithmetic,
  probability, Fourier analysis, and corrected sharp asymptotics},
\newblock manuscript in the \texttt{ProveIt} repository, snapshot 25 August 2026.
\newblock \url{https://github.com/VladimirReshetnikov/ProveIt/blob/main/Analysis/FabiusFunction/docs/Fabius_Function_and_Rvachev_Up/Fabius_Function_and_Rvachev_Up.tex}.

\end{thebibliography}

\end{document}
